\begin{lemma}
Let $\Omega\subseteq\mathbb{C}$ be open, let $S\subseteq\Omega$ meet every
compact subset of $\Omega$ in a finite set, let $j\in\{0,1\}$, let
$g,v\colon\mathbb{C}\to\mathbb{C}$, and let $\beta$ satisfy $0\le\beta<1$.
Assume
\begin{enumerate}
\item[(H1)] $g$ is of class $C^1$ on $\Omega-S$;
\item[(H2)] $\partial_jg=v$ on $\Omega-S$, where $\partial_0$ and $\partial_1$
  are the partial derivatives of \noderef{PartialWord} along the words $[0]$
  and $[1]$;
\item[(H3)] $g\in L^1_{\mathrm{loc}}(\Omega)$ and $v\in L^1_{\mathrm{loc}}(\Omega)$
  in the sense of \noderef{MemLpLocOn};
\item[(H4)] for every $z_0\in S$ there are $M$ and $\rho>0$ with
  $\|g(z)\|\le M\,|z-z_0|^{-\beta}$ whenever $0<|z-z_0|<\rho$.
\end{enumerate}
Then $v$ is the distributional $j$-th partial derivative of $g$ on $\Omega$, in
the sense of \noderef{IsWeakPartialOn}.

This is the step hidden in Section 6's ``Hence
$u\in W^{k+2,q}_{\mathrm{loc}}(\Omega)$'' at paper.tex:534: the derivatives of
the regular representative are classical only off the singular set, and near a
singular point they grow like a negative power of the distance.  The same
statement carries Section 7's classical $p$-harmonic identity off the origin to
the weak form, with $S=\{0\}$ and $\beta=0$ there.
\end{lemma}
\begin{proof}
By \noderef{IsWeakPartialOn} we must show
\begin{equation}\tag{G}
  \int_{\mathbb{C}}g\,\partial_j\varphi\,d\sigma
  =-\int_{\mathbb{C}}v\,\varphi\,d\sigma
\end{equation}
for every test function $\varphi$ for $\Omega$ in the sense of
\noderef{IsTestFn}, that is, every $\varphi$ smooth on $\mathbb{C}$ with
compact support $T=\operatorname{supp}\varphi\subseteq\Omega$.  Fix such a
$\varphi$; write $\|\varphi\|_\infty$ and $\|\partial_j\varphi\|_\infty$ for the
suprema of $|\varphi|$ and $|\partial_j\varphi|$, finite because both are
continuous with compact support.  Throughout $d\sigma$ is plane Lebesgue
measure.

\emph{Step 0: the finitely many exceptional points in the support, and
integrability.}  By hypothesis $S\cap T$ is finite; write
$S\cap T=\{\zeta_1,\dots,\zeta_N\}$ (possibly $N=0$).  Both $g$ and $v$ are
measurable on $\Omega$: they are continuous on the open set $\Omega-S$ --- $g$
by (H1), and $v$ by (H1) together with (H2), being there the continuous
function $\partial_jg$ --- and $S\cap T$ is finite, hence null.  By (H3) and
compactness of $T\subseteq\Omega$,
\begin{equation}\tag{1}
  \int_{T}\|g\|\,d\sigma<\infty,\qquad \int_{T}\|v\|\,d\sigma<\infty ,
\end{equation}
so both integrals in (G) converge absolutely, their integrands vanishing off
$T$.

\emph{Step 1: the cutoffs.}  Fix once and for all a smooth
$\chi\colon\mathbb{R}\to\mathbb{R}$ with
\[
  0\le\chi\le1,\qquad \chi(t)=0\ (t\le1),\qquad \chi(t)=1\ (t\ge2),\qquad
  |\chi'|\le C_0
\]
for an absolute constant $C_0$; such a $\chi$ exists by the standard
construction of smooth bump functions (take $\chi(t)=\int_{-\infty}^{t}\zeta$
with $\zeta\ge0$ smooth, supported in $(1,2)$, of total integral $1$).  By
(H4) choose for each $i\le N$ numbers $M_i$ and $\rho_i>0$ with
\begin{equation}\tag{2}
  \|g(z)\|\le M_i\,|z-\zeta_i|^{-\beta}\qquad(0<|z-\zeta_i|<\rho_i),
\end{equation}
and put $\rho_*=\min_i\rho_i>0$ (with $\rho_*=1$ if $N=0$) and
$M_*=\max_iM_i$.  Fix also
\[
  0<\varepsilon<\tfrac14\min\Bigl\{\rho_*,\ \min_{i\neq i'}|\zeta_i-\zeta_{i'}|\Bigr\},
\]
the second minimum being over distinct pairs and omitted if $N\le1$, and set
\[
  \eta_\varepsilon(z)=\prod_{i=1}^{N}\chi\Bigl(\frac{|z-\zeta_i|}{\varepsilon}\Bigr)
\]
(the empty product being $1$ when $N=0$).  Each factor is smooth on
$\mathbb{C}$: it vanishes identically on the ball $|z-\zeta_i|<\varepsilon$, and
on the open set $z\neq\zeta_i$ the map $z\mapsto|z-\zeta_i|$ is smooth, so the
factor is smooth on each of two open sets covering $\mathbb{C}$.  Hence
$\eta_\varepsilon$ is smooth, and
\begin{equation}\tag{3}
  \eta_\varepsilon=0\ \text{near each }\zeta_i,\qquad
  \eta_\varepsilon=1\ \text{off }\textstyle\bigcup_i\{|z-\zeta_i|<2\varepsilon\},
  \qquad 0\le\eta_\varepsilon\le1 .
\end{equation}
Because $4\varepsilon$ is less than the distance between distinct $\zeta_i$, the
annuli $\varepsilon<|z-\zeta_i|<2\varepsilon$ are pairwise disjoint and at most
one factor of $\eta_\varepsilon$ differs from $1$ at any given point; so by the
product rule, the chain rule and $|\nabla|z-\zeta_i||=1$,
\begin{equation}\tag{4}
  \bigl|\partial_j\eta_\varepsilon(z)\bigr|\le\frac{C_0}{\varepsilon},
  \qquad
  \partial_j\eta_\varepsilon=0\ \text{ unless }
  \varepsilon<|z-\zeta_i|<2\varepsilon\ \text{ for some }i .
\end{equation}

\emph{Step 2: the product has vanishing total derivative.}  Put
\[
  \psi_\varepsilon=g\,\varphi\,\eta_\varepsilon .
\]
We claim $\psi_\varepsilon\in C^1_c(\mathbb{C})$.  Its support is contained in
the compact $T$.  For smoothness, the open sets
\[
  \{z:|z-\zeta_i|<\varepsilon\}\ (i\le N),\qquad \Omega-S,\qquad \mathbb{C}-T
\]
cover $\mathbb{C}$: a point outside all the first ones and outside the third
lies in $T\subseteq\Omega$ and is distinct from every $\zeta_i$, hence lies
outside $S\cap T$; and since it lies in $T$, lying in $S$ would put it in
$S\cap T$, so it lies in $\Omega-S$.  On each of these sets $\psi_\varepsilon$
is $C^1$: on the first ones it is $0$ by (3), on the second it is a product of
the $C^1$ function $g$ with the smooth $\varphi$ and $\eta_\varepsilon$, and on
the third it is $0$ because $\varphi$ is.  Hence $\psi_\varepsilon$ is $C^1$ on
$\mathbb{C}$ with compact support, and therefore
\begin{equation}\tag{5}
  \int_{\mathbb{C}}\partial_j\psi_\varepsilon\,d\sigma=0 .
\end{equation}
Indeed, choose $R$ with $\operatorname{supp}\psi_\varepsilon\subseteq[-R,R]^2$;
for $j=0$, by Fubini and the fundamental theorem of calculus in the first
variable,
\[
  \int_{\mathbb{C}}\partial_x\psi_\varepsilon\,d\sigma
  =\int_{-R}^{R}\Bigl(\int_{-R}^{R}\partial_x\psi_\varepsilon(x,y)\,dx\Bigr)dy
  =\int_{-R}^{R}\bigl(\psi_\varepsilon(R,y)-\psi_\varepsilon(-R,y)\bigr)dy=0,
\]
both boundary values vanishing; the case $j=1$ is the same in the second
variable.  (This is the plane divergence theorem for a $C^1$ field of compact
support, applied to the field whose only nonzero component is the $j$-th.)

\emph{Step 3: expanding (5).}  At every $z\in\Omega-S$ the product rule and
(H2) give
\[
  \partial_j\psi_\varepsilon
  =v\,\varphi\,\eta_\varepsilon+g\,(\partial_j\varphi)\,\eta_\varepsilon
    +g\,\varphi\,(\partial_j\eta_\varepsilon).
\]
At a point $\zeta_i$ both sides vanish, since $\psi_\varepsilon$ is identically
$0$ near $\zeta_i$ and since $\eta_\varepsilon$ and $\partial_j\eta_\varepsilon$
vanish there by (3) and (4); at a point of $S$ outside $T$, and at any point
outside $\Omega$, all terms vanish because $\varphi$ and $\partial_j\varphi$ are
supported in $T\subseteq\Omega$.  So the identity holds everywhere, and all
three terms are integrable by (1), (3) and (4).  Integrating and using (5),
\begin{equation}\tag{6}
  0=\int v\varphi\,\eta_\varepsilon\,d\sigma
   +\int g\,(\partial_j\varphi)\,\eta_\varepsilon\,d\sigma
   +\int g\,\varphi\,(\partial_j\eta_\varepsilon)\,d\sigma .
\end{equation}

\emph{Step 4: the first two terms converge.}  By (3), for every $z$ outside the
finite set $\{\zeta_1,\dots,\zeta_N\}$ we have $\eta_\varepsilon(z)=1$ once
$\varepsilon$ is small enough, so $\eta_\varepsilon\to1$ pointwise off a null
set.  Since
$|v\varphi\eta_\varepsilon|\le\|\varphi\|_\infty\|v\|\mathbf{1}_T$ and
$|g(\partial_j\varphi)\eta_\varepsilon|\le\|\partial_j\varphi\|_\infty\|g\|\mathbf{1}_T$,
both integrable by (1), dominated convergence gives
\[
  \int v\varphi\,\eta_\varepsilon\,d\sigma\to\int v\varphi\,d\sigma,
  \qquad
  \int g\,(\partial_j\varphi)\,\eta_\varepsilon\,d\sigma
  \to\int g\,\partial_j\varphi\,d\sigma .
\]

\emph{Step 5: the cutoff term vanishes --- this is where $\beta<1$ is used.}
By (4) the integrand of the third term of (6) is supported in
$\bigcup_{i\le N}A_{i,\varepsilon}$ with
$A_{i,\varepsilon}=\{z:\varepsilon<|z-\zeta_i|<2\varepsilon\}$, and
$2\varepsilon<\rho_i$, so (2) applies on $A_{i,\varepsilon}$.  Hence, by (4),
\[
  \Bigl|\int g\,\varphi\,(\partial_j\eta_\varepsilon)\,d\sigma\Bigr|
  \le\|\varphi\|_\infty\frac{C_0}{\varepsilon}\sum_{i=1}^{N}
    \int_{A_{i,\varepsilon}}\|g\|\,d\sigma
  \le\|\varphi\|_\infty\frac{C_0M_*}{\varepsilon}\sum_{i=1}^{N}
    \int_{A_{i,\varepsilon}}|z-\zeta_i|^{-\beta}d\sigma .
\]
In polar coordinates about $\zeta_i$,
\[
  \int_{A_{i,\varepsilon}}|z-\zeta_i|^{-\beta}d\sigma
  =2\pi\int_{\varepsilon}^{2\varepsilon}r^{1-\beta}\,dr
  \le2\pi(2\varepsilon)^{1-\beta}\cdot\varepsilon
  =2^{2-\beta}\pi\,\varepsilon^{2-\beta},
\]
using that $r\mapsto r^{1-\beta}$ is nondecreasing, as $1-\beta>0$, so its
maximum on $[\varepsilon,2\varepsilon]$ is $(2\varepsilon)^{1-\beta}$, and that
the interval has length $\varepsilon$.  Therefore
\[
  \Bigl|\int g\,\varphi\,(\partial_j\eta_\varepsilon)\,d\sigma\Bigr|
  \le 2^{2-\beta}\pi C_0M_*N\|\varphi\|_\infty\,\varepsilon^{1-\beta}
  \ \longrightarrow\ 0\qquad(\varepsilon\to0^{+}),
\]
because $1-\beta>0$.  This is the only place the hypothesis $\beta<1$ enters,
and it is sharp: at $\beta=1$ the bound is a constant.  Note that $N$ is fixed
once $\varphi$ is, which is why the local finiteness of $S$ is needed and a
merely discrete $S$ would not do.

\emph{Step 6: conclusion.}  Letting $\varepsilon\to0^{+}$ in (6) and using
Steps 4 and 5,
\[
  0=\int v\varphi\,d\sigma+\int g\,\partial_j\varphi\,d\sigma ,
\]
that is $\int g\,\partial_j\varphi\,d\sigma=-\int v\varphi\,d\sigma$, which is
(G).  As $\varphi$ was an arbitrary test function for $\Omega$, this is the
assertion.
\end{proof}
